\documentclass[11pt]{article}
\usepackage[a4paper,margin=1in]{geometry}
\usepackage{booktabs}
\usepackage{longtable}
\usepackage{hyperref}
\title{Expert Hypothesis Recommendation Paper: p-vs-np}
\date{2026-03-16T22:26:28.893056+00:00}
\begin{document}
\maketitle
\section*{Abstract}
This paper ranks hypotheses using a generative potential function that balances support, coherence, explanatory value, structural leverage, strategy leverage, anti-barrier fit, and testability.
\section*{Generative Function}
\[
Potential(h) = 0.21S + 0.14C + 0.14E + 0.16L + 0.12G + 0.08A + 0.07T + 0.08N
\]
\[
Risk(h) = 0.45R + 0.22C^- + 0.18B + 0.15U
\]
\[
Score(h) = \mathrm{clip}(Potential(h) \cdot (1 - 0.55 Risk(h)) + 0.05V, 0, 1)
\]
\section*{Main Conceptual Bridges}
\begin{longtable}{p{0.18\linewidth}p{0.17\linewidth}p{0.13\linewidth}p{0.20\linewidth}p{0.20\linewidth}}
\toprule
Bridge & Type & Priority & Consequence & Proof Family \\
\midrule
\endhead
bridge.001 & magnification bridge & 0.7437 & NP-not-in-P/poly & invariant-lift-through-intermediate-object \\
bridge.002 & magnification bridge & 0.7319 & NP-not-in-P/poly & invariant-lift-through-intermediate-object \\
bridge.016 & magnification bridge & 0.7108 & exponential-lower-bound & invariant-lift-through-intermediate-object \\
bridge.007 & magnification bridge & 0.6509 & MCSP-hardness-transfer & surrogate-reduction-transfer \\
bridge.010 & magnification bridge & 0.4809 & NP-not-in-P/poly &  \\
bridge.003 & proof-complexity bridge & 0.4719 & TAUT-proof-system-characterization &  \\
\bottomrule
\end{longtable}
\section*{Bridge-First Sketch Rationale}
\subsection*{1. bridge.001}
Move from `st.def.hardness-magnification-phenomenon-programme-leve` on problem `MCSP` in model `nonuniform-circuits` using `A barely-nontrivial lower bound for an explicit problem Q in a weak model` toward `st.thm.oliveira-pich-santhanam-magnification-near-state` in model `nonuniform-circuits` to obtain `major consequences such as NP not in P/poly`.
\begin{itemize}
\item Type: magnification bridge
\item Abstraction: MCSP witness family
\item Carrier: family of instances equipped with a monotone witness measure
\item Typing status: typed
\item Invariant: MCSP complexity profile in nonuniform-circuits
\item Consequence: NP-not-in-P/poly
\item Proof family: invariant-lift-through-intermediate-object
\item Move from `st.def.hardness-magnification-phenomenon-programme-leve` on problem `MCSP` in model `nonuniform-circuits` using `A barely-nontrivial lower bound for an explicit problem Q in a weak model` toward `st.thm.oliveira-pich-santhanam-magnification-near-state` in model `nonuniform-circuits` to obtain `major consequences such as NP not in P/poly`.
\item model shift: nonuniform-circuits / threshold shift: major separation consequence / consequence shift: upgrade `A barely-nontrivial lower bound for an explicit problem Q in a weak model` into `major consequences such as NP not in P/poly`
\item The conjectured object `MCSP witness family` is treated as a `intermediate-complexity-measure` with carrier `family of instances equipped with a monotone witness measure`, playing the role `An explicit family of instances equipped with a monotone witness measure that records enough of `A barely-nontrivial lower bound for an explicit problem Q in a weak model` to test whether the route can reach `NP-not-in-P/poly` while preserving `MCSP complexity profile in nonuniform-circuits`.`.
\item The route-specific missing lemma is: Prove that a hardness-magnification witness family turns `A barely-nontrivial lower bound for an explicit problem Q in a weak model` into `major consequences such as NP not in P/poly` while preserving `MCSP complexity profile in nonuniform-circuits`.; the chosen proof family `invariant-lift-through-intermediate-object` is used because Introduce a conjectured intermediate object or invariant, prove that the source regime induces it, then force the target consequence through that invariant.
\item Bridge priority is supported by effective\_bridge\_score=0.638, raw\_bridge\_score=0.638, conjecture\_priority=0.671, sketch\_score=0.905, hypothesis\_support=0.683, mathematical\_specificity=1.000, evidence\_signal=0.188, falsifiability\_signal=1.000, sketch\_operationalization=1.000, search\_novelty\_signal=1.000.
\end{itemize}
\paragraph{Source Theorem}
A hardness magnification theorem has the form:"A barely-nontrivial lower bound for an explicit problem Q in a weak model implies a major lower bound/separation in a strong model."Typically, Q is a meta-problem (e.g. a Gap-MCSP variant) and"major"means consequences such as NP not in P/poly.
\paragraph{Target Theorem}
In refined regimes, lower bounds for Gap-MCSP slightly beyond known frontiers (e.g. N\textasciicircum{} 1+varepsilon -type circuit lower bounds for suitable promise gaps and input length N=2\textasciicircum{}n ) imply major consequences such as NP not in P/poly.
\paragraph{Mathematical Delta}
\begin{itemize}
\item Model shift: nonuniform-circuits
\item Threshold shift: major separation consequence
\item Consequence shift: upgrade `A barely-nontrivial lower bound for an explicit problem Q in a weak model` into `major consequences such as NP not in P/poly`
\item Missing route lemma: Prove that a hardness-magnification witness family turns `A barely-nontrivial lower bound for an explicit problem Q in a weak model` into `major consequences such as NP not in P/poly` while preserving `MCSP complexity profile in nonuniform-circuits`.
\end{itemize}
\paragraph{Candidate Bridge Object}
Treat `MCSP witness family` as a `intermediate-complexity-measure` attached to the route `st.def.hardness-magnification-phenomenon-programme-leve` -> `st.thm.oliveira-pich-santhanam-magnification-near-state` and typed as a family of instances equipped with a monotone witness measure, playing the role of `An explicit family of instances equipped with a monotone witness measure that records enough of `A barely-nontrivial lower bound for an explicit problem Q in a weak model` to test whether the route can reach `NP-not-in-P/poly` while preserving `MCSP complexity profile in nonuniform-circuits`.`, carrying the source-side content `A barely-nontrivial lower bound for an explicit problem Q in a weak model` across nonuniform-circuits toward `major consequences such as NP not in P/poly` while preserving `MCSP complexity profile in nonuniform-circuits`. Minimal admissibility conditions: admits source-side witness data for `A barely-nontrivial lower bound for an explicit problem Q in a weak model`; preserves the route invariant `MCSP complexity profile in nonuniform-circuits`; forces the target-side payoff `NP-not-in-P/poly`.
\paragraph{Admissibility Conditions}
\begin{itemize}
\item The source-side content `A barely-nontrivial lower bound for an explicit problem Q in a weak model` is available at `st.def.hardness-magnification-phenomenon-programme-leve`.
\item There exists an admissible realization of `MCSP witness family` as a family of instances equipped with a monotone witness measure.
\item The realization preserves `MCSP complexity profile in nonuniform-circuits` across nonuniform-circuits.
\end{itemize}
\paragraph{Lemma Chain}
\begin{itemize}
\item Source-to-object lemma: the source-side content `A barely-nontrivial lower bound for an explicit problem Q in a weak model` at `st.def.hardness-magnification-phenomenon-programme-leve` yields a realization of `MCSP witness family` as a family of instances equipped with a monotone witness measure, satisfying admits source-side witness data for `A barely-nontrivial lower bound for an explicit problem Q in a weak model`; preserves the route invariant `MCSP complexity profile in nonuniform-circuits`; forces the target-side payoff `NP-not-in-P/poly`.
\item Transport lemma: every admissible realization of `MCSP witness family` preserves `MCSP complexity profile in nonuniform-circuits` while crossing nonuniform-circuits.
\item Payoff lemma: if `MCSP witness family` is realized, carries `MCSP complexity profile in nonuniform-circuits`, and discharges the route-specific step from `A barely-nontrivial lower bound for an explicit problem Q in a weak model` to `major consequences such as NP not in P/poly`, then `NP-not-in-P/poly` holds at `st.thm.oliveira-pich-santhanam-magnification-near-state`.
\item Route lemma: Prove that a hardness-magnification witness family turns `A barely-nontrivial lower bound for an explicit problem Q in a weak model` into `major consequences such as NP not in P/poly` while preserving `MCSP complexity profile in nonuniform-circuits`.
\end{itemize}
\paragraph{Proof obligations}
\begin{itemize}
\item Prove the source-to-object claim: the source-side content `A barely-nontrivial lower bound for an explicit problem Q in a weak model` available at `st.def.hardness-magnification-phenomenon-programme-leve` can be organized into `MCSP witness family` as a family of instances equipped with a monotone witness measure.
\item Prove the payoff claim: once `MCSP witness family` carries the route data from `A barely-nontrivial lower bound for an explicit problem Q in a weak model` toward `major consequences such as NP not in P/poly`, it forces `NP-not-in-P/poly` at `st.thm.oliveira-pich-santhanam-magnification-near-state`.
\item Prove the transport claim: every admissible realization of `MCSP witness family` preserves `MCSP complexity profile in nonuniform-circuits` while crossing the route nonuniform-circuits.
\item Discharge the route lemma: Prove that a hardness-magnification witness family turns `A barely-nontrivial lower bound for an explicit problem Q in a weak model` into `major consequences such as NP not in P/poly` while preserving `MCSP complexity profile in nonuniform-circuits`.
\end{itemize}
\paragraph{Failure Diagnosis}
Failure mode: `MCSP witness family` never becomes more than topical overlap because the route-specific lemma remains missing. Prove that a hardness-magnification witness family turns `A barely-nontrivial lower bound for an explicit problem Q in a weak model` into `major consequences such as NP not in P/poly` while preserving `MCSP complexity profile in nonuniform-circuits`.
\paragraph{Negative route}
Allow `MCSP witness family` to ignore the source-side content `A barely-nontrivial lower bound for an explicit problem Q in a weak model` and only preserve topical overlap. bridge-degrades-or-fails The bridge object must absorb the actual source-side mathematical content rather than merely name the antecedent.
\subsection*{2. bridge.002}
Move from `st.def.hardness-magnification-phenomenon-programme-leve` on problem `MCSP` in model `nonuniform-circuits` using `A barely-nontrivial lower bound for an explicit problem Q in a weak model` toward `st.thm.oliveira-santhanam-hardness-magnification-for-na` in model `formula-size` to obtain `major lower bounds for NP in much stronger models (including superpolynomial formula/circuit lower bounds, and in some parameter...`.
\begin{itemize}
\item Type: magnification bridge
\item Abstraction: MCSP witness family
\item Carrier: family of instances equipped with a monotone witness measure
\item Typing status: typed
\item Invariant: MCSP complexity profile in nonuniform-circuits
\item Consequence: NP-not-in-P/poly
\item Proof family: invariant-lift-through-intermediate-object
\item Move from `st.def.hardness-magnification-phenomenon-programme-leve` on problem `MCSP` in model `nonuniform-circuits` using `A barely-nontrivial lower bound for an explicit problem Q in a weak model` toward `st.thm.oliveira-santhanam-hardness-magnification-for-na` in model `formula-size` to obtain `major lower bounds for NP in much stronger models (including superpolynomial formula/circuit lower bounds, and in some parameter...`.
\item model shift: nonuniform-circuits -> formula-size / threshold shift: major separation consequence / consequence shift: upgrade `A barely-nontrivial lower bound for an explicit problem Q in a weak model` into `major lower bounds for NP in much stronger models (including superpolynomial formula/circuit lower bounds, and in some parameter...`
\item The conjectured object `MCSP witness family` is treated as a `intermediate-complexity-measure` with carrier `family of instances equipped with a monotone witness measure`, playing the role `An explicit family of instances equipped with a monotone witness measure that records enough of `A barely-nontrivial lower bound for an explicit problem Q in a weak model` to test whether the route can reach `NP-not-in-P/poly` while preserving `MCSP complexity profile in nonuniform-circuits`.`.
\item The route-specific missing lemma is: Prove that a hardness-magnification witness family preserves `MCSP complexity profile in nonuniform-circuits` while transporting the route from the nonuniform-circuits regime to the formula-size regime strongly enough to derive `major lower bounds for NP in much stronger models (including superpolynomial formula/circuit lower bounds, and in some parameter...`.; the chosen proof family `invariant-lift-through-intermediate-object` is used because Introduce a conjectured intermediate object or invariant, prove that the source regime induces it, then force the target consequence through that invariant.
\item Bridge priority is supported by effective\_bridge\_score=0.585, raw\_bridge\_score=0.585, conjecture\_priority=0.651, sketch\_score=0.897, hypothesis\_support=0.683, mathematical\_specificity=1.000, evidence\_signal=0.188, falsifiability\_signal=1.000, sketch\_operationalization=1.000, search\_novelty\_signal=1.000.
\end{itemize}
\paragraph{Source Theorem}
A hardness magnification theorem has the form:"A barely-nontrivial lower bound for an explicit problem Q in a weak model implies a major lower bound/separation in a strong model."Typically, Q is a meta-problem (e.g. a Gap-MCSP variant) and"major"means consequences such as NP not in P/poly.
\paragraph{Target Theorem}
There exist explicit"meta"problems (including MCSP variants) such that proving certain emph barely-superlinear lower bounds against restricted models (e.g. formulas) would imply major lower bounds for NP in much stronger models (including superpolynomial formula/circuit lower bounds, and in some parameter regimes consequences as strong as NP not in P/poly ).
\paragraph{Mathematical Delta}
\begin{itemize}
\item Model shift: nonuniform-circuits -> formula-size
\item Threshold shift: major separation consequence
\item Consequence shift: upgrade `A barely-nontrivial lower bound for an explicit problem Q in a weak model` into `major lower bounds for NP in much stronger models (including superpolynomial formula/circuit lower bounds, and in some parameter...`
\item Missing route lemma: Prove that a hardness-magnification witness family preserves `MCSP complexity profile in nonuniform-circuits` while transporting the route from the nonuniform-circuits regime to the formula-size regime strongly enough to derive `major lower bounds for NP in much stronger models (including superpolynomial formula/circuit lower bounds, and in some parameter...`.
\end{itemize}
\paragraph{Candidate Bridge Object}
Treat `MCSP witness family` as a `intermediate-complexity-measure` attached to the route `st.def.hardness-magnification-phenomenon-programme-leve` -> `st.thm.oliveira-santhanam-hardness-magnification-for-na` and typed as a family of instances equipped with a monotone witness measure, playing the role of `An explicit family of instances equipped with a monotone witness measure that records enough of `A barely-nontrivial lower bound for an explicit problem Q in a weak model` to test whether the route can reach `NP-not-in-P/poly` while preserving `MCSP complexity profile in nonuniform-circuits`.`, carrying the source-side content `A barely-nontrivial lower bound for an explicit problem Q in a weak model` across nonuniform-circuits -> formula-size toward `major lower bounds for NP in much stronger models (including superpolynomial formula/circuit lower bounds, and in some parameter...` while preserving `MCSP complexity profile in nonuniform-circuits`. Minimal admissibility conditions: admits source-side witness data for `A barely-nontrivial lower bound for an explicit problem Q in a weak model`; preserves the route invariant `MCSP complexity profile in nonuniform-circuits`; forces the target-side payoff `NP-not-in-P/poly`.
\paragraph{Admissibility Conditions}
\begin{itemize}
\item The source-side content `A barely-nontrivial lower bound for an explicit problem Q in a weak model` is available at `st.def.hardness-magnification-phenomenon-programme-leve`.
\item There exists an admissible realization of `MCSP witness family` as a family of instances equipped with a monotone witness measure.
\item The realization preserves `MCSP complexity profile in nonuniform-circuits` across nonuniform-circuits -> formula-size.
\end{itemize}
\paragraph{Lemma Chain}
\begin{itemize}
\item Source-to-object lemma: the source-side content `A barely-nontrivial lower bound for an explicit problem Q in a weak model` at `st.def.hardness-magnification-phenomenon-programme-leve` yields a realization of `MCSP witness family` as a family of instances equipped with a monotone witness measure, satisfying admits source-side witness data for `A barely-nontrivial lower bound for an explicit problem Q in a weak model`; preserves the route invariant `MCSP complexity profile in nonuniform-circuits`; forces the target-side payoff `NP-not-in-P/poly`.
\item Transport lemma: every admissible realization of `MCSP witness family` preserves `MCSP complexity profile in nonuniform-circuits` while crossing nonuniform-circuits -> formula-size.
\item Payoff lemma: if `MCSP witness family` is realized, carries `MCSP complexity profile in nonuniform-circuits`, and discharges the route-specific step from `A barely-nontrivial lower bound for an explicit problem Q in a weak model` to `major lower bounds for NP in much stronger models (including superpolynomial formula/circuit lower bounds, and in some parameter...`, then `NP-not-in-P/poly` holds at `st.thm.oliveira-santhanam-hardness-magnification-for-na`.
\item Route lemma: Prove that a hardness-magnification witness family preserves `MCSP complexity profile in nonuniform-circuits` while transporting the route from the nonuniform-circuits regime to the formula-size regime strongly enough to derive `major lower bounds for NP in much stronger models (including superpolynomial formula/circuit lower bounds, and in some parameter...`.
\end{itemize}
\paragraph{Proof obligations}
\begin{itemize}
\item Prove the source-to-object claim: the source-side content `A barely-nontrivial lower bound for an explicit problem Q in a weak model` available at `st.def.hardness-magnification-phenomenon-programme-leve` can be organized into `MCSP witness family` as a family of instances equipped with a monotone witness measure.
\item Prove the payoff claim: once `MCSP witness family` carries the route data from `A barely-nontrivial lower bound for an explicit problem Q in a weak model` toward `major lower bounds for NP in much stronger models (including superpolynomial formula/circuit lower bounds, and in some parameter...`, it forces `NP-not-in-P/poly` at `st.thm.oliveira-santhanam-hardness-magnification-for-na`.
\item Prove the transport claim: every admissible realization of `MCSP witness family` preserves `MCSP complexity profile in nonuniform-circuits` while crossing the route nonuniform-circuits -> formula-size.
\item Discharge the route lemma: Prove that a hardness-magnification witness family preserves `MCSP complexity profile in nonuniform-circuits` while transporting the route from the nonuniform-circuits regime to the formula-size regime strongly enough to derive `major lower bounds for NP in much stronger models (including superpolynomial formula/circuit lower bounds, and in some parameter...`.
\end{itemize}
\paragraph{Failure Diagnosis}
Failure mode: `MCSP witness family` never becomes more than topical overlap because the route-specific lemma remains missing. Prove that a hardness-magnification witness family preserves `MCSP complexity profile in nonuniform-circuits` while transporting the route from the nonuniform-circuits regime to the formula-size regime strongly enough to derive `major lower bounds for NP in much stronger models (including superpolynomial formula/circuit lower bounds, and in some parameter...`.
\paragraph{Negative route}
Allow `MCSP witness family` to ignore the source-side content `A barely-nontrivial lower bound for an explicit problem Q in a weak model` and only preserve topical overlap. bridge-degrades-or-fails The bridge object must absorb the actual source-side mathematical content rather than merely name the antecedent.
\subsection*{3. bridge.016}
Move from `st.def.hardness-magnification-phenomenon-programme-leve` on problem `MCSP` in model `nonuniform-circuits` using `A barely-nontrivial lower bound for an explicit problem Q in a weak model` toward `st.thm.atserias-m-uller-simple-general-magnification-fr` to obtain `There exist general-purpose magnification theorems (not problem-specific) showing that certain mild lower bounds for...`.
\begin{itemize}
\item Type: magnification bridge
\item Abstraction: MCSP witness family
\item Carrier: family of instances equipped with a monotone witness measure
\item Typing status: typed
\item Invariant: exponential-lower-bound witness complexity
\item Consequence: exponential-lower-bound
\item Proof family: invariant-lift-through-intermediate-object
\item Move from `st.def.hardness-magnification-phenomenon-programme-leve` on problem `MCSP` in model `nonuniform-circuits` using `A barely-nontrivial lower bound for an explicit problem Q in a weak model` toward `st.thm.atserias-m-uller-simple-general-magnification-fr` to obtain `There exist general-purpose magnification theorems (not problem-specific) showing that certain mild lower bounds for...`.
\item model shift: nonuniform-circuits / threshold shift: major separation consequence -> exponential lower bound / consequence shift: upgrade `A barely-nontrivial lower bound for an explicit problem Q in a weak model` into `There exist general-purpose magnification theorems (not problem-specific) showing that certain mild lower bounds for...`
\item The conjectured object `MCSP witness family` is treated as a `intermediate-complexity-measure` with carrier `family of instances equipped with a monotone witness measure`, playing the role `An explicit family of instances equipped with a monotone witness measure that records enough of `A barely-nontrivial lower bound for an explicit problem Q in a weak model` to test whether the route can reach `exponential-lower-bound` while preserving `exponential-lower-bound witness complexity`.`.
\item The route-specific missing lemma is: Prove that a hardness-magnification witness family upgrades the route from major separation consequence to exponential lower bound without losing `exponential-lower-bound witness complexity`.; the chosen proof family `invariant-lift-through-intermediate-object` is used because Introduce a conjectured intermediate object or invariant, prove that the source regime induces it, then force the target consequence through that invariant.
\item Bridge priority is supported by effective\_bridge\_score=0.495, raw\_bridge\_score=0.495, conjecture\_priority=0.609, sketch\_score=0.890, hypothesis\_support=0.659, mathematical\_specificity=1.000, evidence\_signal=0.200, falsifiability\_signal=1.000, sketch\_operationalization=1.000, search\_novelty\_signal=1.000.
\end{itemize}
\paragraph{Source Theorem}
A hardness magnification theorem has the form:"A barely-nontrivial lower bound for an explicit problem Q in a weak model implies a major lower bound/separation in a strong model."Typically, Q is a meta-problem (e.g. a Gap-MCSP variant) and"major"means consequences such as NP not in P/poly.
\paragraph{Target Theorem}
There exist general-purpose magnification theorems (not problem-specific) showing that certain mild lower bounds for distinguishers/sparse-code objects can be amplified into stronger circuit/formula lower bounds, achieving magnification thresholds below previously known lower bounds in some settings.
\paragraph{Mathematical Delta}
\begin{itemize}
\item Model shift: nonuniform-circuits
\item Threshold shift: major separation consequence -> exponential lower bound
\item Consequence shift: upgrade `A barely-nontrivial lower bound for an explicit problem Q in a weak model` into `There exist general-purpose magnification theorems (not problem-specific) showing that certain mild lower bounds for...`
\item Missing route lemma: Prove that a hardness-magnification witness family upgrades the route from major separation consequence to exponential lower bound without losing `exponential-lower-bound witness complexity`.
\end{itemize}
\paragraph{Candidate Bridge Object}
Treat `MCSP witness family` as a `intermediate-complexity-measure` attached to the route `st.def.hardness-magnification-phenomenon-programme-leve` -> `st.thm.atserias-m-uller-simple-general-magnification-fr` and typed as a family of instances equipped with a monotone witness measure, playing the role of `An explicit family of instances equipped with a monotone witness measure that records enough of `A barely-nontrivial lower bound for an explicit problem Q in a weak model` to test whether the route can reach `exponential-lower-bound` while preserving `exponential-lower-bound witness complexity`.`, carrying the source-side content `A barely-nontrivial lower bound for an explicit problem Q in a weak model` across nonuniform-circuits toward `There exist general-purpose magnification theorems (not problem-specific) showing that certain mild lower bounds for...` while preserving `exponential-lower-bound witness complexity`. Minimal admissibility conditions: admits source-side witness data for `A barely-nontrivial lower bound for an explicit problem Q in a weak model`; preserves the route invariant `exponential-lower-bound witness complexity`; aligns the model shift (nonuniform-circuits) with an explicit transport map; forces the target-side payoff `exponential-lower-bound`.
\paragraph{Admissibility Conditions}
\begin{itemize}
\item The source-side content `A barely-nontrivial lower bound for an explicit problem Q in a weak model` is available at `st.def.hardness-magnification-phenomenon-programme-leve`.
\item There exists an admissible realization of `MCSP witness family` as a family of instances equipped with a monotone witness measure.
\item The realization preserves `exponential-lower-bound witness complexity` across nonuniform-circuits.
\end{itemize}
\paragraph{Lemma Chain}
\begin{itemize}
\item Source-to-object lemma: the source-side content `A barely-nontrivial lower bound for an explicit problem Q in a weak model` at `st.def.hardness-magnification-phenomenon-programme-leve` yields a realization of `MCSP witness family` as a family of instances equipped with a monotone witness measure, satisfying admits source-side witness data for `A barely-nontrivial lower bound for an explicit problem Q in a weak model`; preserves the route invariant `exponential-lower-bound witness complexity`; aligns the model shift (nonuniform-circuits) with an explicit transport map.
\item Transport lemma: every admissible realization of `MCSP witness family` preserves `exponential-lower-bound witness complexity` while crossing nonuniform-circuits.
\item Payoff lemma: if `MCSP witness family` is realized, carries `exponential-lower-bound witness complexity`, and discharges the route-specific step from `A barely-nontrivial lower bound for an explicit problem Q in a weak model` to `There exist general-purpose magnification theorems (not problem-specific) showing that certain mild lower bounds for...`, then `exponential-lower-bound` holds at `st.thm.atserias-m-uller-simple-general-magnification-fr`.
\item Route lemma: Prove that a hardness-magnification witness family upgrades the route from major separation consequence to exponential lower bound without losing `exponential-lower-bound witness complexity`.
\end{itemize}
\paragraph{Proof obligations}
\begin{itemize}
\item Prove the source-to-object claim: the source-side content `A barely-nontrivial lower bound for an explicit problem Q in a weak model` available at `st.def.hardness-magnification-phenomenon-programme-leve` can be organized into `MCSP witness family` as a family of instances equipped with a monotone witness measure.
\item Prove the payoff claim: once `MCSP witness family` carries the route data from `A barely-nontrivial lower bound for an explicit problem Q in a weak model` toward `There exist general-purpose magnification theorems (not problem-specific) showing that certain mild lower bounds for...`, it forces `exponential-lower-bound` at `st.thm.atserias-m-uller-simple-general-magnification-fr`.
\item Prove the transport claim: every admissible realization of `MCSP witness family` preserves `exponential-lower-bound witness complexity` while crossing the route nonuniform-circuits.
\item Exhibit the model-alignment claim: `MCSP witness family` gives a faithful translation across nonuniform-circuits.
\end{itemize}
\paragraph{Failure Diagnosis}
Failure mode: `MCSP witness family` never becomes more than topical overlap because the route-specific lemma remains missing. Prove that a hardness-magnification witness family upgrades the route from major separation consequence to exponential lower bound without losing `exponential-lower-bound witness complexity`.
\paragraph{Negative route}
Allow `MCSP witness family` to ignore the source-side content `A barely-nontrivial lower bound for an explicit problem Q in a weak model` and only preserve topical overlap. bridge-degrades-or-fails The bridge object must absorb the actual source-side mathematical content rather than merely name the antecedent.
\subsection*{4. bridge.007}
Move from `st.thm.karp-many-natural-problems-are-np-complete-karp1` using `NP-complete under polynomial-time many-one reductions` toward `st.thm.murray-williams-non-np-hardness-of-mcsp-under-lo` on problem `MCSP` to obtain `is not NP-complete. More broadly, the paper develops evidence and conditional consequences suggesting NP-hardness of MCSP under various...`.
\begin{itemize}
\item Type: magnification bridge
\item Abstraction: unspecified witness family
\item Carrier: family of instances equipped with a monotone witness measure
\item Typing status: typed
\item Invariant: consequence schema NP-complete-classification
\item Consequence: MCSP-hardness-transfer
\item Proof family: surrogate-reduction-transfer
\item Move from `st.thm.karp-many-natural-problems-are-np-complete-karp1` using `NP-complete under polynomial-time many-one reductions` toward `st.thm.murray-williams-non-np-hardness-of-mcsp-under-lo` on problem `MCSP` to obtain `is not NP-complete. More broadly, the paper develops evidence and conditional consequences suggesting NP-hardness of MCSP under various...`.
\item model shift: model shift not yet typed / threshold shift: threshold regime not yet typed / consequence shift: upgrade `NP-complete under polynomial-time many-one reductions` into `is not NP-complete. More broadly, the paper develops evidence and conditional consequences suggesting NP-hardness of MCSP under various...`
\item The conjectured object `unspecified witness family` is treated as a `intermediate-complexity-measure` with carrier `family of instances equipped with a monotone witness measure`, playing the role `An explicit family of instances equipped with a monotone witness measure that records enough of `NP-complete under polynomial-time many-one reductions` to test whether the route can reach `MCSP-hardness-transfer` while preserving `consequence schema NP-complete-classification`.`.
\item The route-specific missing lemma is: Prove that a hardness-magnification witness family turns `NP-complete under polynomial-time many-one reductions` into `is not NP-complete. More broadly, the paper develops evidence and conditional consequences suggesting NP-hardness of MCSP under various...` while preserving `consequence schema NP-complete-classification`.; the chosen proof family `surrogate-reduction-transfer` is used because Define a surrogate object or promise problem, reduce the source statement into it, and transfer the required consequence through the surrogate interface.
\item Bridge priority is supported by effective\_bridge\_score=0.509, raw\_bridge\_score=0.509, conjecture\_priority=0.685, sketch\_score=0.871, hypothesis\_support=0.643, mathematical\_specificity=0.700, evidence\_signal=0.250, falsifiability\_signal=0.750, sketch\_operationalization=1.000, search\_novelty\_signal=1.000.
\end{itemize}
\paragraph{Source Theorem}
There exist many problems L in combinatorial optimisation/graph theory such that L is NP-complete under polynomial-time many-one reductions.
\paragraph{Target Theorem}
Under"local"polynomial-time reductions (where each output bit is computable with very small time), MCSP is not NP-complete. More broadly, the paper develops evidence and conditional consequences suggesting NP-hardness of MCSP under various reductions would entail unexpected class separations.
\paragraph{Mathematical Delta}
\begin{itemize}
\item Model shift: model shift not yet typed
\item Threshold shift: threshold regime not yet typed
\item Consequence shift: upgrade `NP-complete under polynomial-time many-one reductions` into `is not NP-complete. More broadly, the paper develops evidence and conditional consequences suggesting NP-hardness of MCSP under various...`
\item Missing route lemma: Prove that a hardness-magnification witness family turns `NP-complete under polynomial-time many-one reductions` into `is not NP-complete. More broadly, the paper develops evidence and conditional consequences suggesting NP-hardness of MCSP under various...` while preserving `consequence schema NP-complete-classification`.
\end{itemize}
\paragraph{Candidate Bridge Object}
Treat `unspecified witness family` as a `intermediate-complexity-measure` attached to the route `st.thm.karp-many-natural-problems-are-np-complete-karp1` -> `st.thm.murray-williams-non-np-hardness-of-mcsp-under-lo` and typed as a family of instances equipped with a monotone witness measure, playing the role of `An explicit family of instances equipped with a monotone witness measure that records enough of `NP-complete under polynomial-time many-one reductions` to test whether the route can reach `MCSP-hardness-transfer` while preserving `consequence schema NP-complete-classification`.`, carrying the source-side content `NP-complete under polynomial-time many-one reductions` across model shift not yet typed toward `is not NP-complete. More broadly, the paper develops evidence and conditional consequences suggesting NP-hardness of MCSP under various...` while preserving `consequence schema NP-complete-classification`. Minimal admissibility conditions: admits source-side witness data for `NP-complete under polynomial-time many-one reductions`; preserves the route invariant `consequence schema NP-complete-classification`; provides an explicit transfer law or reduction mechanism on the route; forces the target-side payoff `MCSP-hardness-transfer`.
\paragraph{Admissibility Conditions}
\begin{itemize}
\item The source-side content `NP-complete under polynomial-time many-one reductions` is available at `st.thm.karp-many-natural-problems-are-np-complete-karp1`.
\item There exists an admissible realization of `unspecified witness family` as a family of instances equipped with a monotone witness measure.
\item The realization preserves `consequence schema NP-complete-classification` across model shift not yet typed.
\end{itemize}
\paragraph{Lemma Chain}
\begin{itemize}
\item Source-to-object lemma: the source-side content `NP-complete under polynomial-time many-one reductions` at `st.thm.karp-many-natural-problems-are-np-complete-karp1` yields a realization of `unspecified witness family` as a family of instances equipped with a monotone witness measure, satisfying admits source-side witness data for `NP-complete under polynomial-time many-one reductions`; preserves the route invariant `consequence schema NP-complete-classification`; provides an explicit transfer law or reduction mechanism on the route.
\item Transport lemma: every admissible realization of `unspecified witness family` preserves `consequence schema NP-complete-classification` while crossing model shift not yet typed.
\item Payoff lemma: if `unspecified witness family` is realized, carries `consequence schema NP-complete-classification`, and discharges the route-specific step from `NP-complete under polynomial-time many-one reductions` to `is not NP-complete. More broadly, the paper develops evidence and conditional consequences suggesting NP-hardness of MCSP under various...`, then `MCSP-hardness-transfer` holds at `st.thm.murray-williams-non-np-hardness-of-mcsp-under-lo`.
\item Reduction lemma: the source-side content `NP-complete under polynomial-time many-one reductions` reduces to the surrogate object encoded by `unspecified witness family` without weakening the path to `MCSP-hardness-transfer`.
\end{itemize}
\paragraph{Proof obligations}
\begin{itemize}
\item Prove the source-to-object claim: the source-side content `NP-complete under polynomial-time many-one reductions` available at `st.thm.karp-many-natural-problems-are-np-complete-karp1` can be organized into `unspecified witness family` as a family of instances equipped with a monotone witness measure.
\item Prove the payoff claim: once `unspecified witness family` carries the route data from `NP-complete under polynomial-time many-one reductions` toward `is not NP-complete. More broadly, the paper develops evidence and conditional consequences suggesting NP-hardness of MCSP under various...`, it forces `MCSP-hardness-transfer` at `st.thm.murray-williams-non-np-hardness-of-mcsp-under-lo`.
\item Prove the transport claim: every admissible realization of `unspecified witness family` preserves `consequence schema NP-complete-classification` while crossing the route model shift not yet typed.
\item Construct the transfer-law claim: an explicit reduction or transport map sends `NP-complete under polynomial-time many-one reductions` into `unspecified witness family` without losing the route to `MCSP-hardness-transfer`.
\end{itemize}
\paragraph{Failure Diagnosis}
Failure mode: `unspecified witness family` never becomes more than topical overlap because the route-specific lemma remains missing. Prove that a hardness-magnification witness family turns `NP-complete under polynomial-time many-one reductions` into `is not NP-complete. More broadly, the paper develops evidence and conditional consequences suggesting NP-hardness of MCSP under various...` while preserving `consequence schema NP-complete-classification`.
\paragraph{Negative route}
Allow `unspecified witness family` to ignore the source-side content `NP-complete under polynomial-time many-one reductions` and only preserve topical overlap. bridge-degrades-or-fails The bridge object must absorb the actual source-side mathematical content rather than merely name the antecedent.
\subsection*{5. bridge.010}
Move from `st.thm.oliveira-pich-santhanam-magnification-near-state` on problem `MCSP` in model `nonuniform-circuits` using `In refined regimes, lower bounds for Gap-MCSP slightly beyond known frontiers (e.g. N\textasciicircum{} 1+varepsilon -type circuit lower bounds for...` toward `st.thm.oliveira-santhanam-hardness-magnification-for-na` in model `formula-size` to obtain `major lower bounds for NP in much stronger models (including superpolynomial formula/circuit lower bounds, and in some parameter...`.
\begin{itemize}
\item Type: magnification bridge
\item Abstraction: MCSP witness family
\item Carrier: family of instances equipped with a monotone witness measure
\item Typing status: typed
\item Invariant: MCSP complexity profile in nonuniform-circuits
\item Consequence: NP-not-in-P/poly
\item Proof family: 
\item Move from `st.thm.oliveira-pich-santhanam-magnification-near-state` on problem `MCSP` in model `nonuniform-circuits` using `In refined regimes, lower bounds for Gap-MCSP slightly beyond known frontiers (e.g. N\textasciicircum{} 1+varepsilon -type circuit lower bounds for...` toward `st.thm.oliveira-santhanam-hardness-magnification-for-na` in model `formula-size` to obtain `major lower bounds for NP in much stronger models (including superpolynomial formula/circuit lower bounds, and in some parameter...`.
\item model shift: nonuniform-circuits -> formula-size / threshold shift: major separation consequence / consequence shift: upgrade `In refined regimes, lower bounds for Gap-MCSP slightly beyond known frontiers (e.g. N\textasciicircum{} 1+varepsilon -type circuit lower bounds for...` into `major lower bounds for NP in much stronger models (including superpolynomial formula/circuit lower bounds, and in some parameter...`
\item The conjectured object `MCSP witness family` is treated as a `intermediate-complexity-measure` with carrier `family of instances equipped with a monotone witness measure`, playing the role `An explicit family of instances equipped with a monotone witness measure that records enough of `In refined regimes, lower bounds for Gap-MCSP slightly beyond known frontiers (e.g. N\textasciicircum{} 1+varepsilon -type circuit lower bounds for...` to test whether the route can reach `NP-not-in-P/poly` while preserving `MCSP complexity profile in nonuniform-circuits`.`.
\item The route-specific missing lemma is: Prove that a MCSP witness family preserves `MCSP complexity profile in nonuniform-circuits` while transporting the route from the nonuniform-circuits regime to the formula-size regime strongly enough to derive `major lower bounds for NP in much stronger models (including superpolynomial formula/circuit lower bounds, and in some parameter...`.; the chosen proof family `untyped family` is used because it supplies the most concrete source-to-target transport mechanism currently available
\item Bridge priority is supported by effective\_bridge\_score=0.558, raw\_bridge\_score=0.558, conjecture\_priority=0.641, sketch\_score=0.000, hypothesis\_support=0.713, mathematical\_specificity=0.880, evidence\_signal=0.250, falsifiability\_signal=0.250, sketch\_operationalization=0.000, search\_novelty\_signal=1.000.
\end{itemize}
\paragraph{Source Theorem}
In refined regimes, lower bounds for Gap-MCSP slightly beyond known frontiers (e.g. N\textasciicircum{} 1+varepsilon -type circuit lower bounds for suitable promise gaps and input length N=2\textasciicircum{}n ) imply major consequences such as NP not in P/poly.
\paragraph{Target Theorem}
There exist explicit"meta"problems (including MCSP variants) such that proving certain emph barely-superlinear lower bounds against restricted models (e.g. formulas) would imply major lower bounds for NP in much stronger models (including superpolynomial formula/circuit lower bounds, and in some parameter regimes consequences as strong as NP not in P/poly ).
\paragraph{Mathematical Delta}
\begin{itemize}
\item Model shift: nonuniform-circuits -> formula-size
\item Threshold shift: major separation consequence
\item Consequence shift: upgrade `In refined regimes, lower bounds for Gap-MCSP slightly beyond known frontiers (e.g. N\textasciicircum{} 1+varepsilon -type circuit lower bounds for...` into `major lower bounds for NP in much stronger models (including superpolynomial formula/circuit lower bounds, and in some parameter...`
\item Missing route lemma: Prove that a MCSP witness family preserves `MCSP complexity profile in nonuniform-circuits` while transporting the route from the nonuniform-circuits regime to the formula-size regime strongly enough to derive `major lower bounds for NP in much stronger models (including superpolynomial formula/circuit lower bounds, and in some parameter...`.
\end{itemize}
\paragraph{Candidate Bridge Object}
An adequate `intermediate-complexity-measure` bridge concept for `bridge.010` is a `family of instances equipped with a monotone witness measure` named `MCSP witness family` that absorbs `MCSP-near-threshold-nonuniform-circuit-lower-bound`, preserves `MCSP complexity profile in nonuniform-circuits`, and is strong enough to force `NP-not-in-P/poly`.
\paragraph{Admissibility Conditions}
\begin{itemize}
\item The source-side content `In refined regimes, lower bounds for Gap-MCSP slightly beyond known frontiers (e.g. N\textasciicircum{} 1+varepsilon -type circuit lower bounds for...` is available in st.thm.oliveira-pich-santhanam-magnification-near-state (problems MCSP; models nonuniform-circuits; claims In refined regimes, lower bounds for Gap-MCSP slightly beyond known frontiers (e.g. N\textasciicircum{} 1+varepsilon -type circuit lower bounds for..., MCSP nonuniform circuits lower bound).
\item There exists a `MCSP witness family` object, typed as a family of instances equipped with a monotone witness measure.
\item The object preserves `MCSP complexity profile in nonuniform-circuits` across the source-to-target move nonuniform-circuits -> formula-size.
\end{itemize}
\paragraph{Lemma Chain}
\begin{itemize}
\item No explicit candidate lemmas surfaced yet.
\end{itemize}
\paragraph{Proof obligations}
\begin{itemize}
\item No explicit obligations surfaced yet.
\end{itemize}
\paragraph{Failure Diagnosis}
Failure mode: `MCSP witness family` never becomes more than topical overlap because the route-specific lemma remains missing. Prove that a MCSP witness family preserves `MCSP complexity profile in nonuniform-circuits` while transporting the route from the nonuniform-circuits regime to the formula-size regime strongly enough to derive `major lower bounds for NP in much stronger models (including superpolynomial formula/circuit lower bounds, and in some parameter...`.
\paragraph{Negative route}
No explicit negative route surfaced yet.
\subsection*{6. bridge.003}
Move from `st.def.cook-reckhow-propositional-proof-system` on problem `TAUT` in model `Cook-Reckhow-proof-system` using `A (Cook--Reckhow) propositional proof system is a polynomial-time computable surjection f:bits\textasciicircum{}*twoheadrightarrow TAUT, where TAUT is...` toward `st.thm.cook-reckhow-proof-systems-and-cookreckhow1979` in model `Cook-Reckhow-proof-system` to obtain `There exists a polynomially bounded propositional proof system for TAUT if and only if NP=coNP`.
\begin{itemize}
\item Type: proof-complexity bridge
\item Abstraction: TAUT proof progress measure
\item Carrier: monotone potential function over proofs or certificates
\item Typing status: typed
\item Invariant: TAUT complexity profile in Cook-Reckhow-proof-system
\item Consequence: TAUT-proof-system-characterization
\item Proof family: 
\item Move from `st.def.cook-reckhow-propositional-proof-system` on problem `TAUT` in model `Cook-Reckhow-proof-system` using `A (Cook--Reckhow) propositional proof system is a polynomial-time computable surjection f:bits\textasciicircum{}*twoheadrightarrow TAUT, where TAUT is...` toward `st.thm.cook-reckhow-proof-systems-and-cookreckhow1979` in model `Cook-Reckhow-proof-system` to obtain `There exists a polynomially bounded propositional proof system for TAUT if and only if NP=coNP`.
\item model shift: Cook-Reckhow-proof-system / threshold shift: threshold regime not yet typed / consequence shift: upgrade `A (Cook--Reckhow) propositional proof system is a polynomial-time computable surjection f:bits\textasciicircum{}*twoheadrightarrow TAUT, where TAUT is...` into `There exists a polynomially bounded propositional proof system for TAUT if and only if NP=coNP`
\item The conjectured object `TAUT proof progress measure` is treated as a `proof-potential-measure` with carrier `monotone potential function over proofs or certificates`, playing the role `A monotone potential function over proofs or certificates witnessing whether the move from `A (Cook--Reckhow) propositional proof system is a polynomial-time computable surjection f:bits\textasciicircum{}*twoheadrightarrow TAUT, where TAUT is...` toward `TAUT-proof-system-characterization` is structurally achievable.`.
\item The route-specific missing lemma is: Prove that a proof certificate family or monotone proof measure turns `A (Cook--Reckhow) propositional proof system is a polynomial-time computable surjection f:bits\textasciicircum{}*twoheadrightarrow TAUT, where TAUT is...` into `There exists a polynomially bounded propositional proof system for TAUT if and only if NP=coNP` while preserving `TAUT complexity profile in Cook-Reckhow-proof-system`.; the chosen proof family `untyped family` is used because it supplies the most concrete source-to-target transport mechanism currently available
\item Bridge priority is supported by effective\_bridge\_score=0.554, raw\_bridge\_score=0.554, conjecture\_priority=0.632, sketch\_score=0.000, hypothesis\_support=0.641, mathematical\_specificity=0.880, evidence\_signal=0.250, falsifiability\_signal=0.250, sketch\_operationalization=0.000, search\_novelty\_signal=1.000.
\end{itemize}
\paragraph{Source Theorem}
A (Cook--Reckhow) propositional proof system is a polynomial-time computable surjection f:bits\textasciicircum{}*twoheadrightarrow TAUT, where TAUT is the set of propositional tautologies. A string pi is an f -proof of formula varphi if f(pi)=varphi.
\paragraph{Target Theorem}
There exists a polynomially bounded propositional proof system for TAUT if and only if NP=coNP.
\paragraph{Mathematical Delta}
\begin{itemize}
\item Model shift: Cook-Reckhow-proof-system
\item Threshold shift: threshold regime not yet typed
\item Consequence shift: upgrade `A (Cook--Reckhow) propositional proof system is a polynomial-time computable surjection f:bits\textasciicircum{}*twoheadrightarrow TAUT, where TAUT is...` into `There exists a polynomially bounded propositional proof system for TAUT if and only if NP=coNP`
\item Missing route lemma: Prove that a proof certificate family or monotone proof measure turns `A (Cook--Reckhow) propositional proof system is a polynomial-time computable surjection f:bits\textasciicircum{}*twoheadrightarrow TAUT, where TAUT is...` into `There exists a polynomially bounded propositional proof system for TAUT if and only if NP=coNP` while preserving `TAUT complexity profile in Cook-Reckhow-proof-system`.
\end{itemize}
\paragraph{Candidate Bridge Object}
An adequate `proof-potential-measure` bridge concept for `bridge.003` is a `monotone potential function over proofs or certificates` named `TAUT proof progress measure` that absorbs `proof-system-characterization`, preserves `TAUT complexity profile in Cook-Reckhow-proof-system`, and is strong enough to force `TAUT-proof-system-characterization`.
\paragraph{Admissibility Conditions}
\begin{itemize}
\item The source-side content `A (Cook--Reckhow) propositional proof system is a polynomial-time computable surjection f:bits\textasciicircum{}*twoheadrightarrow TAUT, where TAUT is...` is available in st.def.cook-reckhow-propositional-proof-system (problems TAUT; models Cook-Reckhow-proof-system; claims A (Cook--Reckhow) propositional proof system is a polynomial-time computable surjection f:bits\textasciicircum{}*twoheadrightarrow TAUT, where TAUT is..., proof system characterization).
\item There exists a `TAUT proof progress measure` object, typed as a monotone potential function over proofs or certificates.
\item The object preserves `TAUT complexity profile in Cook-Reckhow-proof-system` across the source-to-target move Cook-Reckhow-proof-system.
\end{itemize}
\paragraph{Lemma Chain}
\begin{itemize}
\item No explicit candidate lemmas surfaced yet.
\end{itemize}
\paragraph{Proof obligations}
\begin{itemize}
\item No explicit obligations surfaced yet.
\end{itemize}
\paragraph{Failure Diagnosis}
Failure mode: `TAUT proof progress measure` never becomes more than topical overlap because the route-specific lemma remains missing. Prove that a proof certificate family or monotone proof measure turns `A (Cook--Reckhow) propositional proof system is a polynomial-time computable surjection f:bits\textasciicircum{}*twoheadrightarrow TAUT, where TAUT is...` into `There exists a polynomially bounded propositional proof system for TAUT if and only if NP=coNP` while preserving `TAUT complexity profile in Cook-Reckhow-proof-system`.
\paragraph{Negative route}
No explicit negative route surfaced yet.
\section*{Top Recommendations}
\begin{longtable}{p{0.28\linewidth}p{0.20\linewidth}p{0.13\linewidth}p{0.13\linewidth}p{0.13\linewidth}}
\toprule
Hypothesis & Tier & Score & Potential & Risk \\
\midrule
\endhead
phyp.p-vs-np.033 & Tier-1: Immediate Expert Focus & 0.8243 & 0.8032 & 0.0654 \\
phyp.p-vs-np.001 & Tier-1: Immediate Expert Focus & 0.7932 & 0.7752 & 0.0751 \\
phyp.p-vs-np.002 & Tier-1: Immediate Expert Focus & 0.7387 & 0.7189 & 0.0764 \\
phyp.p-vs-np.010 & Tier-2: High-Upside Program & 0.7127 & 0.6904 & 0.0730 \\
phyp.p-vs-np.042 & Tier-2: High-Upside Program & 0.6671 & 0.6382 & 0.0603 \\
phyp.p-vs-np.035 & Tier-2: High-Upside Program & 0.6633 & 0.6903 & 0.0711 \\
phyp.p-vs-np.044 & Tier-2: High-Upside Program & 0.6632 & 0.6963 & 0.0864 \\
phyp.p-vs-np.034 & Tier-2: High-Upside Program & 0.6568 & 0.6309 & 0.0696 \\
\bottomrule
\end{longtable}
\section*{Detailed Justification}
\subsection*{1. phyp.p-vs-np.033}
\begin{itemize}
\item Tier: Tier-1: Immediate Expert Focus
\item Kind: definition-introduction
\item Target: st.thm.oliveira-pich-santhanam-magnification-near-state
\item Support: st.def.hardness-magnification-phenomenon-programme-leve
\item Evidence balance: support=0.859, refute=0.000, confidence=0.804.
\item Global-theory effect: coherence\_delta=+0.0169, coherence\_signal=0.585, explanatory\_signal=0.539.
\item Structural leverage: entry=0.998, intervention=0.961, pathway=1.000, combined=0.985.
\item Strategy leverage: linked\_strategies=2, strategy\_signal=1.000, concept\_fertility=0.822.
\item Barrier/readiness: barrier\_scope=none, anti\_barrier\_alignment=0.600, testability=1.000, revision\_priority=1.000.
\item Risk decomposition: refutation=0.000, negative\_coherence=0.000, barrier\_pressure=0.201, uncertainty=0.196; aggregated\_risk=0.065.
\end{itemize}
\paragraph{Recommended Expert Actions}
\begin{itemize}
\item Run focused deduction test on `st.def.hardness-magnification-phenomenon-programme-leve -> st.thm.oliveira-pich-santhanam-magnification-near-state` under the current inferential commitment.
\item Instantiate definition-introduction assumption connecting st.def.hardness-magnification-phenomenon-programme-leve -> st.thm.oliveira-pich-santhanam-magnification-near-state in a scratch branch.
\item State the preserved invariant `MCSP complexity profile in nonuniform-circuits` explicitly and test it on both sides.
\item Follow evaluator action: promote-or-retest-with-deduction.
\end{itemize}
\subsection*{2. phyp.p-vs-np.001}
\begin{itemize}
\item Tier: Tier-1: Immediate Expert Focus
\item Kind: bridge-link
\item Target: st.thm.oliveira-pich-santhanam-magnification-near-state
\item Support: st.def.hardness-magnification-phenomenon-programme-leve
\item Evidence balance: support=0.809, refute=0.000, confidence=0.740.
\item Global-theory effect: coherence\_delta=+0.0171, coherence\_signal=0.585, explanatory\_signal=0.451.
\item Structural leverage: entry=0.998, intervention=0.961, pathway=1.000, combined=0.985.
\item Strategy leverage: linked\_strategies=2, strategy\_signal=1.000, concept\_fertility=0.822.
\item Barrier/readiness: barrier\_scope=none, anti\_barrier\_alignment=0.600, testability=1.000, revision\_priority=1.000.
\item Risk decomposition: refutation=0.000, negative\_coherence=0.000, barrier\_pressure=0.201, uncertainty=0.260; aggregated\_risk=0.075.
\end{itemize}
\paragraph{Recommended Expert Actions}
\begin{itemize}
\item Run focused deduction test on `st.def.hardness-magnification-phenomenon-programme-leve -> st.thm.oliveira-pich-santhanam-magnification-near-state` under the current inferential commitment.
\item Add tentative dependency st.thm.oliveira-pich-santhanam-magnification-near-state depends\_on st.def.hardness-magnification-phenomenon-programme-leve in a scratch KB copy.
\item Check whether `MCSP complexity profile in nonuniform-circuits` is preserved after instantiating st.def.hardness-magnification-phenomenon-programme-leve->st.thm.oliveira-pich-santhanam-magnification-near-state.
\item Follow evaluator action: promote-or-retest-with-deduction.
\end{itemize}
\subsection*{3. phyp.p-vs-np.002}
\begin{itemize}
\item Tier: Tier-1: Immediate Expert Focus
\item Kind: bridge-link
\item Target: st.thm.oliveira-santhanam-hardness-magnification-for-na
\item Support: st.def.hardness-magnification-phenomenon-programme-leve
\item Evidence balance: support=0.803, refute=0.000, confidence=0.732.
\item Global-theory effect: coherence\_delta=+0.0078, coherence\_signal=0.539, explanatory\_signal=0.421.
\item Structural leverage: entry=0.998, intervention=1.000, pathway=1.000, combined=0.999.
\item Strategy leverage: linked\_strategies=1, strategy\_signal=0.610, concept\_fertility=0.831.
\item Barrier/readiness: barrier\_scope=none, anti\_barrier\_alignment=0.600, testability=1.000, revision\_priority=1.000.
\item Risk decomposition: refutation=0.000, negative\_coherence=0.000, barrier\_pressure=0.201, uncertainty=0.268; aggregated\_risk=0.076.
\end{itemize}
\paragraph{Recommended Expert Actions}
\begin{itemize}
\item Run focused deduction test on `st.def.hardness-magnification-phenomenon-programme-leve -> st.thm.oliveira-santhanam-hardness-magnification-for-na` under the current inferential commitment.
\item Add tentative dependency st.thm.oliveira-santhanam-hardness-magnification-for-na depends\_on st.def.hardness-magnification-phenomenon-programme-leve in a scratch KB copy.
\item Check whether `MCSP complexity profile in nonuniform-circuits` is preserved after instantiating st.def.hardness-magnification-phenomenon-programme-leve->st.thm.oliveira-santhanam-hardness-magnification-for-na.
\item Follow evaluator action: promote-or-retest-with-deduction.
\end{itemize}
\subsection*{4. phyp.p-vs-np.010}
\begin{itemize}
\item Tier: Tier-2: High-Upside Program
\item Kind: bridge-link
\item Target: st.thm.oliveira-santhanam-hardness-magnification-for-na
\item Support: st.thm.oliveira-pich-santhanam-magnification-near-state
\item Evidence balance: support=0.821, refute=0.000, confidence=0.754.
\item Global-theory effect: coherence\_delta=+0.0075, coherence\_signal=0.538, explanatory\_signal=0.408.
\item Structural leverage: entry=0.965, intervention=1.000, pathway=0.982, combined=0.981.
\item Strategy leverage: linked\_strategies=1, strategy\_signal=0.390, concept\_fertility=0.814.
\item Barrier/readiness: barrier\_scope=none, anti\_barrier\_alignment=0.600, testability=1.000, revision\_priority=1.000.
\item Risk decomposition: refutation=0.000, negative\_coherence=0.000, barrier\_pressure=0.201, uncertainty=0.246; aggregated\_risk=0.073.
\end{itemize}
\paragraph{Recommended Expert Actions}
\begin{itemize}
\item Run focused deduction test on `st.thm.oliveira-pich-santhanam-magnification-near-state -> st.thm.oliveira-santhanam-hardness-magnification-for-na` under the current inferential commitment.
\item Add tentative dependency st.thm.oliveira-santhanam-hardness-magnification-for-na depends\_on st.thm.oliveira-pich-santhanam-magnification-near-state in a scratch KB copy.
\item Check whether `MCSP complexity profile in nonuniform-circuits` is preserved after instantiating st.thm.oliveira-pich-santhanam-magnification-near-state->st.thm.oliveira-santhanam-hardness-magnification-for-na.
\item Follow evaluator action: promote-or-retest-with-deduction.
\end{itemize}
\subsection*{5. phyp.p-vs-np.042}
\begin{itemize}
\item Tier: Tier-2: High-Upside Program
\item Kind: barrier-escape-principle
\item Target: st.thm.oliveira-santhanam-hardness-magnification-for-na
\item Support: st.thm.oliveira-pich-santhanam-magnification-near-state
\item Evidence balance: support=0.885, refute=0.000, confidence=0.839.
\item Global-theory effect: coherence\_delta=+0.0074, coherence\_signal=0.537, explanatory\_signal=0.519.
\item Structural leverage: entry=0.965, intervention=1.000, pathway=0.980, combined=0.980.
\item Strategy leverage: linked\_strategies=0, strategy\_signal=0.000, concept\_fertility=0.000.
\item Barrier/readiness: barrier\_scope=none, anti\_barrier\_alignment=0.600, testability=1.000, revision\_priority=1.000.
\item Risk decomposition: refutation=0.000, negative\_coherence=0.000, barrier\_pressure=0.201, uncertainty=0.161; aggregated\_risk=0.060.
\end{itemize}
\paragraph{Recommended Expert Actions}
\begin{itemize}
\item Run focused deduction test on `st.thm.oliveira-pich-santhanam-magnification-near-state -> st.thm.oliveira-santhanam-hardness-magnification-for-na` under the current inferential commitment.
\item Instantiate barrier-escape-principle assumption connecting st.thm.oliveira-pich-santhanam-magnification-near-state -> st.thm.oliveira-santhanam-hardness-magnification-for-na in a scratch branch.
\item State the preserved invariant `MCSP complexity profile in nonuniform-circuits` explicitly and test it on both sides.
\item Follow evaluator action: promote-or-retest-with-deduction.
\end{itemize}
\subsection*{6. phyp.p-vs-np.035}
\begin{itemize}
\item Tier: Tier-2: High-Upside Program
\item Kind: definition-introduction
\item Target: st.thm.cook-reckhow-proof-systems-and-cookreckhow1979
\item Support: st.def.cook-reckhow-propositional-proof-system
\item Evidence balance: support=0.830, refute=0.000, confidence=0.767.
\item Global-theory effect: coherence\_delta=+0.0048, coherence\_signal=0.524, explanatory\_signal=0.570.
\item Structural leverage: entry=0.926, intervention=0.727, pathway=0.548, combined=0.780.
\item Strategy leverage: linked\_strategies=1, strategy\_signal=0.414, concept\_fertility=0.830.
\item Barrier/readiness: barrier\_scope=none, anti\_barrier\_alignment=0.600, testability=1.000, revision\_priority=0.000.
\item Risk decomposition: refutation=0.000, negative\_coherence=0.000, barrier\_pressure=0.201, uncertainty=0.233; aggregated\_risk=0.071.
\end{itemize}
\paragraph{Recommended Expert Actions}
\begin{itemize}
\item Run focused deduction test on `st.def.cook-reckhow-propositional-proof-system -> st.thm.cook-reckhow-proof-systems-and-cookreckhow1979` under the current inferential commitment.
\item Instantiate definition-introduction assumption connecting st.def.cook-reckhow-propositional-proof-system -> st.thm.cook-reckhow-proof-systems-and-cookreckhow1979 in a scratch branch.
\item State the preserved invariant `TAUT complexity profile in Cook-Reckhow-proof-system` explicitly and test it on both sides.
\item Follow evaluator action: promote-or-retest-with-deduction.
\end{itemize}
\subsection*{7. phyp.p-vs-np.044}
\begin{itemize}
\item Tier: Tier-2: High-Upside Program
\item Kind: definition-introduction
\item Target: st.thm.karp-many-natural-problems-are-np-complete-karp1
\item Support: st.def.np-completeness
\item Evidence balance: support=0.750, refute=0.000, confidence=0.665.
\item Global-theory effect: coherence\_delta=+0.0047, coherence\_signal=0.523, explanatory\_signal=0.521.
\item Structural leverage: entry=0.905, intervention=0.708, pathway=0.678, combined=0.791.
\item Strategy leverage: linked\_strategies=1, strategy\_signal=0.605, concept\_fertility=0.943.
\item Barrier/readiness: barrier\_scope=none, anti\_barrier\_alignment=0.600, testability=1.000, revision\_priority=0.000.
\item Risk decomposition: refutation=0.000, negative\_coherence=0.000, barrier\_pressure=0.201, uncertainty=0.335; aggregated\_risk=0.086.
\end{itemize}
\paragraph{Recommended Expert Actions}
\begin{itemize}
\item Run focused deduction test on `st.def.np-completeness -> st.thm.karp-many-natural-problems-are-np-complete-karp1` under the current inferential commitment.
\item Instantiate definition-introduction assumption connecting st.def.np-completeness -> st.thm.karp-many-natural-problems-are-np-complete-karp1 in a scratch branch.
\item State the preserved invariant `consequence schema NP-complete-classification` explicitly and test it on both sides.
\item Follow evaluator action: promote-or-retest-with-deduction.
\end{itemize}
\subsection*{8. phyp.p-vs-np.034}
\begin{itemize}
\item Tier: Tier-2: High-Upside Program
\item Kind: definition-introduction
\item Target: st.thm.oliveira-santhanam-hardness-magnification-for-na
\item Support: st.def.hardness-magnification-phenomenon-programme-leve
\item Evidence balance: support=0.838, refute=0.000, confidence=0.777.
\item Global-theory effect: coherence\_delta=+0.0076, coherence\_signal=0.538, explanatory\_signal=0.510.
\item Structural leverage: entry=0.998, intervention=1.000, pathway=1.000, combined=0.999.
\item Strategy leverage: linked\_strategies=0, strategy\_signal=0.000, concept\_fertility=0.000.
\item Barrier/readiness: barrier\_scope=none, anti\_barrier\_alignment=0.600, testability=1.000, revision\_priority=1.000.
\item Risk decomposition: refutation=0.000, negative\_coherence=0.000, barrier\_pressure=0.201, uncertainty=0.223; aggregated\_risk=0.070.
\end{itemize}
\paragraph{Recommended Expert Actions}
\begin{itemize}
\item Run focused deduction test on `st.def.hardness-magnification-phenomenon-programme-leve -> st.thm.oliveira-santhanam-hardness-magnification-for-na` under the current inferential commitment.
\item Instantiate definition-introduction assumption connecting st.def.hardness-magnification-phenomenon-programme-leve -> st.thm.oliveira-santhanam-hardness-magnification-for-na in a scratch branch.
\item State the preserved invariant `MCSP complexity profile in nonuniform-circuits` explicitly and test it on both sides.
\item Follow evaluator action: promote-or-retest-with-deduction.
\end{itemize}
\end{document}
